\documentclass[10pt,a4paper]{amsart}
\usepackage[margin=19mm]{geometry}
\usepackage{amsmath,amssymb,mathtools}
\usepackage[scr=rsfs,cal=euler]{mathalfa}
\usepackage{xfrac}
\usepackage[unicode,hidelinks,bookmarks=false]{hyperref}
\newcommand{\ie}{i.e.\ }
\newcommand{\cf}{cf.\ }
\newcommand{\resp}{respectively\ }
\newcommand{\Subsection}{\subsection}
\providecommand{\nobreakdash}{\nobreak\hbox{-}}
\setlength{\emergencystretch}{2em}
\multlinegap0pt
\usepackage{bm}
\usepackage{bbm}
\usepackage{dsfont}
\usepackage{xspace}
\usepackage{cancel}
\usepackage{mathtools}
\usepackage[shortlabels]{enumitem}
\usepackage{amsthm}
\makeatletter
\@ifundefined{theorem}{\newtheorem{theorem}{Theorem}}{}
\@ifundefined{lemma}{\newtheorem{lemma}[theorem]{Lemma}}{}
\makeatother
\newcommand{\E}{\mathbb{E}}
\newcommand{\Ortho}{\operatorname{Ortho}}
\newcommand{\R}{\mathbb{R}}
\newcommand{\abs}[1]{\left\lvert #1\right\rvert}
\newcommand{\ip}[2]{\left\langle #1,#2\right\rangle}
\newcommand{\norm}[1]{\left\lVert #1\right\rVert}
\title[Sharp First-Order Lower Bounds under $\alpha$-Polyak-Lojasie]{Sharp First-Order Lower Bounds under $\alpha$-Polyak-Lojasiewicz Conditions}
\author{Saeed Masiha, Negar Kiyavash,, Patrick Thiran}
\date{Source: arXiv:2606.28278v3 (CC BY 4.0)}
\begin{document}
\maketitle

\noindent\emph{Source and licence.} Sharp First-Order Lower Bounds under $\alpha$-Polyak-Lojasiewicz Conditions. Version arXiv:2606.28278v3 (CC BY 4.0), distributed under the Creative Commons Attribution 4.0 International licence, as recorded in the arXiv licence metadata for this exact version and shown on the abstract page https://arxiv.org/abs/2606.28278v3. Full rights notice: licenses/arxiv-2606.28278-CC-BY-4.0.txt.\par\smallskip

\noindent\textit{Editorial scope.} Counted unit: the Lemma whose source label is \texttt{lem:endpoint-global-pl} in the article (source0.tex lines 5697--5708, source label \texttt{lem:endpoint-global-pl}), delivered together with the article's own complete original proof of it (source0.tex lines 5710--6015, one \texttt{proof} environment of 306 lines). No mathematics is written, repaired or re-derived: statement and proof are the article's own text line for line. Required context retained uncounted: lem:Gamma-basic (lines 3785--3793); lem:Gamma-outside-box (lines 3818--3824); lem:Gamma-box-strong (lines 3861--3864); lem:radial-clipping (lines 4407--4416); eq:hatF-def (lines 4533--4536); lem:clipped-rotated (lines 4551--4584). Every definition, display and label of those lines is retained unchanged. Omitted from this package and not redistributed: 1--3784, 3794--3817, 3825--3860, 3865--4406, 4417--4532, 4537--4550, 4585--5696, 5709--5709, 6016--6433. Nothing inside a retained excerpt is omitted or altered: every retained body is delivered exactly as the article's lines are written, so no omission receipt of any kind accompanies this package. Citations: the retained excerpts carry no citation command at all, so no bibliography entry of the article is transcribed and no reference is redistributed. Reference adaptation: none was needed and none was made; every reference inside the delivered unit resolves inside it, so no unresolved reference or citation is produced. Printed slips kept literal: no printed word, symbol, hypothesis or number of the counted statement or of its proof was altered. Layout and class: the article's own class is \texttt{article} with packages including geometry, amsmath, amssymb, amsthm, mathtools, fontenc, bm, natbib, authblk, hyperref, xcolor, enumitem, array, booktabs; the wrapper is the lane's pinned \texttt{amsart} editorial wrapper at 10pt on A4, which restates the theorem-like environments the retained text needs and adds the article's own macro definitions that the retained text needs (\texttt{\textbackslash E}, \texttt{\textbackslash Ortho}, \texttt{\textbackslash R}, \texttt{\textbackslash abs}, \texttt{\textbackslash ip}, \texttt{\textbackslash norm}), with extra wrapper packages (bm, bbm, dsfont, xspace, cancel, mathtools, enumitem). The delivered numbering is the unit's own. Assets: no retained span contains a figure environment, a graphics-inclusion command, a PDF-page inclusion, a table or a TikZ picture; no image file is copied into this package, the package declares no asset dependency and no image file is redistributed with it. Licence: version arXiv:2606.28278v3 of this article is distributed under the Creative Commons Attribution 4.0 International licence, as recorded in the arXiv licence metadata for this exact version and shown on the abstract page https://arxiv.org/abs/2606.28278v3. A full rights notice is delivered at licenses/arxiv-2606.28278-CC-BY-4.0.txt. Characters: every retained excerpt is pure ASCII, so the delivered engine, XeLaTeX with its default Latin Modern text font, reports no missing character.
\begin{lemma}[Basic properties of $\Gamma_N$]
\label{lem:Gamma-basic}
For every $N\ge 1$:
\begin{enumerate}[label=\arabic*.,leftmargin=1.5em]
    \item $\Gamma_N$ is coercive, hence attains its minimum.
    \item $\Gamma_N(0)=F_N(0)\le 3N$.
    \item $S_0(\Gamma_N)\subseteq S_0(F_N)$.
\end{enumerate}
\end{lemma}

\begin{lemma}[Gradient bounded away from zero outside the box for $\Gamma_N$]
\label{lem:Gamma-outside-box}
For every $N\ge 1$ and every $x\in \R^N\setminus \mathcal B_N$,
\[
\norm{\nabla \Gamma_N(x)}\ge \frac{1}{20}.
\]
\end{lemma}

\begin{lemma}[Local strong convexity on the box for $\Gamma_N$]
\label{lem:Gamma-box-strong}
$\Gamma_N$ is $(1/2+\eta)$-strongly convex on $\mathcal B_N$.
\end{lemma}

\begin{lemma}[Global regularity of the radial clipping map]
\label{lem:radial-clipping}
For every $m\ge 1$ and every Euclidean dimension \(d\), the Jacobian $J\rho_m$ is globally
$(11/R_m)$-Lipschitz on \(\R^d\):
\[
\norm{J\rho_m(x)-J\rho_m(y)}_{\mathrm{op}}\le \frac{11}{R_m}\norm{x-y}
\qquad
(\forall x,y\in\R^d).
\]
\end{lemma}

\begin{equation}
\label{eq:hatF-def}
\widehat F_{m,U}(x):=F_{2m}(U^\top \rho_m(x))+\frac{\eta}{2}\norm{x}^2,
\end{equation}

\begin{lemma}[Properties of the clipped, rotated, regularized instance]
\label{lem:clipped-rotated}
There exists a numerical constant $c_0>0$ such that the following hold for every $m\ge m_0$, every $q\in(0,1]$, every $d\ge d_{m,q}$, and every $U\in \Ortho(d,2m)$:
\begin{enumerate}[label=\arabic*.,leftmargin=1.5em]
    \item $\widehat g_{m,q,U}$ is an unbiased gradient estimator for $\widehat F_{m,U}$.
    \item
    \[
    \E\bigl[\norm{\widehat g_{m,q,U}(x,z)-\nabla \widehat F_{m,U}(x)}^2\mid x\bigr]\le \frac{4}{q}.
    \]
    \item $\nabla \widehat F_{m,U}$ is globally $L_{\mathrm{clip}}$-Lipschitz for a universal numerical constant $L_{\mathrm{clip}}$.
    \item $S_0(\widehat F_{m,U})\subseteq B(0,R_m)$.
    Hence, by~\eqref{eq:hatF-def}, on $S_0(\widehat F_{m,U})$,
    \[
    \widehat F_{m,U}(x)=F_{2m}(U^\top x)+\frac{\eta}{2}\norm{x}^2,
    \]
    and
    \[
    \widehat F_{m,U}(x)-\widehat F_{m,U,\star}\le 2400m\norm{\nabla \widehat F_{m,U}(x)}^2
    \qquad
    (x\in S_0(\widehat F_{m,U})).
    \]
    \item If $U$ is drawn uniformly from $\Ortho(d,2m)$, $A$ is any randomized algorithm in the gradient-only stochastic
    first-order model of this paper, and
    \[
    t\le c_0\frac{m}{q},
    \]
    then
    \[
    \E_{U,\xi}\bigl[\widehat F_{m,U}(x_t^A)-\widehat F_{m,U,\star}\bigr]\ge m,
    \]
    where \(x_t^A\) is the point generated when \(A\) is run with stochastic-gradient oracle
    \(\widehat g_{m,q,U}\), and the expectation is over $U$, the oracle randomness, and the internal randomness of $A$.
\end{enumerate}
\end{lemma}

\begin{lemma}[Global regularity and P\L{} geometry of the endpoint instance]
\label{lem:endpoint-global-pl}
For every $m\ge m_0$, every $d\ge2m$, and every $U\in\Ortho(d,2m)$, the gradient of
$\widehat F_{m,U}$ is globally $113$-Lipschitz and
\[
\widehat F_{m,U}(x)-\widehat F_{m,U,\star}
\le
28{,}778m\norm{\nabla\widehat F_{m,U}(x)}^2
\qquad
(\forall x\in\R^d).
\]
\end{lemma}

\begin{proof}
Set $N:=2m$, recall that $\eta=1/20$, and note that
$R_m=16\sqrt m=8\sqrt2\,\sqrt N$. We first sharpen the scalar estimates used by the endpoint construction. We then
prove the smoothness and global P\L{} bounds separately.

\emph{Step 1: global scalar and chain estimates.}
Write
\[
F_N(x)=H_N(x)+\sum_{i=1}^N B(x_i),
\qquad
H_N(x):=\varphi(x_1)+\sum_{i=1}^{N-1}\Theta(x_i)\varphi(x_{i+1}).
\]
Because $0\le\Theta,\varphi\le1$, the $N$ summands in $H_N$ give $0\le H_N\le N$. Each coordinate of
$\nabla H_N$ is the sum of at most one term bounded by $\sup|\varphi'|<1$ and one term bounded by
$\sup|\Theta'|\sup\varphi\le3$. Hence
\begin{equation}
\label{eq:endpoint-H-bounds}
0\le H_N(x)\le N,
\qquad
\norm{\nabla H_N(x)}\le4\sqrt N.
\end{equation}

We also need three estimates for $B$. We verify them on every piece of its definition. If $t\le-1/2$, then
$B(t)=2+(t+1/2)^2/2$ and $B'(t)=t+1/2$, so
\[
B(t)-\left(\frac12t^2+2\right)
=\frac12t+\frac18\le0,
\qquad
\abs{B'(t)-t}=\frac12,
\]
and
\[
tB'(t)-\left(\frac12t^2-\frac{13}{2}\right)
=\frac12\left(t+\frac12\right)^2+\frac{51}{8}\ge0.
\]
If $-1/2<t\le1/2$, then $B(t)=2\le t^2/2+2$, $\abs{B'(t)-t}=\abs{t}\le1/2$, and
$tB'(t)=0\ge t^2/2-13/2$. If
$1/2<t<1$, then $0\le B(t)\le2$ and $B'(t)=-2\Theta'(t)$. Since $0<t<1$ and
$0\le\Theta'(t)\le3$,
\[
B(t)\le2\le\frac12t^2+2,
\qquad
\abs{B'(t)-t}\le7,
\qquad
tB'(t)\ge-6\ge\frac12t^2-\frac{13}{2}.
\]
Finally, if $t\ge1$, then $B(t)=(t-1)^2/2$ and $B'(t)=t-1$, whence
\[
B(t)\le\frac12t^2+2,
\qquad
\abs{B'(t)-t}=1,
\]
and
\[
tB'(t)-\left(\frac12t^2-\frac{13}{2}\right)
=\frac12(t-1)^2+6\ge0.
\]
Thus, globally,
\begin{equation}
\label{eq:endpoint-B-bounds}
B(t)\le\frac12t^2+2,
\qquad
\abs{B'(t)-t}\le7,
\qquad
tB'(t)\ge\frac12t^2-\frac{13}{2}.
\end{equation}

For $r:=\norm{x}$, equations~\eqref{eq:endpoint-H-bounds}--\eqref{eq:endpoint-B-bounds} imply
\begin{equation}
\label{eq:endpoint-Gamma-global-growth}
\Gamma_N(x)
\le
\left(\frac12+\frac{\eta}{2}\right)r^2+3N
=\frac{21}{40}r^2+3N
\end{equation}
and
\begin{equation}
\label{eq:endpoint-Gamma-radial}
\ip{\nabla\Gamma_N(x)}{x}
\ge
\left(\frac12+\eta\right)r^2-4\sqrt N\,r-\frac{13}{2}N
=\frac{11}{20}r^2-4\sqrt N\,r-\frac{13}{2}N.
\end{equation}

\emph{Step 2: the sharper smoothness bound.}
For any $y\in\R^N$ with $\norm{y}\le2R_m$, equations~\eqref{eq:endpoint-H-bounds}--\eqref{eq:endpoint-B-bounds}
give
\[
\norm{\nabla F_N(y)}
\le
\norm{\nabla H_N(y)}+\norm{B'(y)}
\le
\norm{y}+11\sqrt N
\le
3R_m,
\]
where the last inequality uses $\norm{y}\le2R_m$ and $11\sqrt N\le R_m$. For arbitrary $x,x'\in\R^d$, set
$y:=U^\top\rho_m(x)$ and $y':=U^\top\rho_m(x')$. The clipping map is $1$-Lipschitz,
$\sup_x\norm{J\rho_m(x)}_{\mathrm{op}}\le1$, and Lemma~\ref{lem:radial-clipping} gives
$\mathrm{Lip}(J\rho_m)\le11/R_m$. Therefore
\begin{align*}
&\norm{J\rho_m(x)^\top U\nabla F_N(y)-J\rho_m(x')^\top U\nabla F_N(y')}\\
&\qquad\le
79\norm{x-x'}+\frac{11}{R_m}(3R_m)\norm{x-x'}
=112\norm{x-x'}.
\end{align*}
Adding the gradient $\eta x$ of the quadratic regularizer shows that $\nabla\widehat F_{m,U}$ is globally
$(112+\eta)$-Lipschitz, and hence globally $113$-Lipschitz.

\emph{Step 3: a global P\L{} inequality for $\Gamma_N$.}
Let $x^\star$ be a global minimizer of $\Gamma_N$. It exists by Lemma~\ref{lem:Gamma-basic}(1), and
Lemma~\ref{lem:Gamma-outside-box} implies $x^\star\in\mathcal B_N$ because
$\nabla\Gamma_N(x^\star)=0$. If $x\in\mathcal B_N$, Lemma~\ref{lem:Gamma-box-strong} and the standard gradient
bound for strongly convex functions give
\[
\Gamma_N(x)-\Gamma_{N,\star}\le\norm{\nabla\Gamma_N(x)}^2.
\]

Suppose next that $x\notin\mathcal B_N$ and $r\le9\sqrt N$. The four cases in the proof of
Lemma~\ref{lem:Gamma-outside-box} actually give the sharper component bounds
$1/8$, $9/160$, $3/40$, and $1/8$, respectively. Hence
$\norm{\nabla\Gamma_N(x)}\ge9/160$. Meanwhile,
\eqref{eq:endpoint-Gamma-global-growth} and $\Gamma_{N,\star}\ge0$ give
\[
\Gamma_N(x)-\Gamma_{N,\star}
\le\frac{1821}{40}N
\le14{,}389N\norm{\nabla\Gamma_N(x)}^2.
\]

It remains to consider $x\notin\mathcal B_N$ and $r>9\sqrt N$. The difference between the right-hand side of
\eqref{eq:endpoint-Gamma-radial} and $r^2/40$ is
\[
\frac{21}{40}r^2-4\sqrt N\,r-\frac{13}{2}N.
\]
This expression is increasing for $r\ge9\sqrt N$ and equals $N/40>0$ at $r=9\sqrt N$. Therefore
$\ip{\nabla\Gamma_N(x)}{x}\ge r^2/40$, and Cauchy--Schwarz yields
$\norm{\nabla\Gamma_N(x)}\ge r/40$. Using~\eqref{eq:endpoint-Gamma-global-growth} and
$3N\le r^2/27$ once more,
\[
\Gamma_N(x)-\Gamma_{N,\star}
\le
\left(\frac{21}{40}+\frac{1}{27}\right)r^2
\le900\norm{\nabla\Gamma_N(x)}^2.
\]
Combining the three regions gives the global estimate
\begin{equation}
\label{eq:endpoint-Gamma-global-PL}
\Gamma_N(x)-\Gamma_{N,\star}
\le
C_\Gamma N\norm{\nabla\Gamma_N(x)}^2
\qquad(\forall x\in\R^N),
\quad
C_\Gamma:=14{,}389.
\end{equation}

\emph{Step 4: points inside the clipping ball.}
Assume $\norm{x}\le R_m$, so $\rho_m(x)=x$, and define
\[
y:=U^\top x,
\qquad
w:=x-Uy.
\]
The vectors $Uy$ and $w$ are orthogonal. Step~4 in the proof of Lemma~\ref{lem:clipped-rotated} establishes
$\widehat F_{m,U,\star}=\Gamma_{N,\star}$; it follows by minimizing $\Gamma_N$ and observing that its minimizer lies
strictly inside the clipping ball. Consequently,
\begin{equation}
\label{eq:endpoint-inside-decomposition}
\widehat F_{m,U}(x)
=\Gamma_N(y)+\frac{\eta}{2}\norm{w}^2,
\qquad
\norm{\nabla\widehat F_{m,U}(x)}^2
=\norm{\nabla\Gamma_N(y)}^2+\eta^2\norm{w}^2.
\end{equation}
Since $N=2m$, equations~\eqref{eq:endpoint-Gamma-global-PL}--\eqref{eq:endpoint-inside-decomposition} give
\[
\widehat F_{m,U}(x)-\widehat F_{m,U,\star}
\le
28{,}778m\norm{\nabla\Gamma_N(y)}^2+10\eta^2\norm{w}^2
\le
28{,}778m\norm{\nabla\widehat F_{m,U}(x)}^2.
\]

\emph{Step 5: points outside the clipping ball.}
Let $r:=\norm{x}>R_m$ and introduce
\begin{align*}
e&:=\frac{x}{r},
&
h&:=h_m(r),
&
a&:=\frac{h}{r},
&
z&:=U^\top e,\\
y&:=hz,
&
g&:=\nabla F_N(y),
&
\beta&:=\ip{g}{z}.
\end{align*}
The radial and tangential eigenvalues of $J\rho_m(x)$ are $h_m'(r)$ and $a$, respectively. Since
$\ip{Ug}{e}=\beta$, the two orthogonal gradient components are
\begin{equation}
\label{eq:endpoint-exterior-gradient}
\nabla\widehat F_{m,U}(x)
=
\bigl(h_m'(r)\beta+\eta r\bigr)e
+a\bigl(Ug-\beta e\bigr),
\end{equation}
and hence
\begin{equation}
\label{eq:endpoint-exterior-gradient-norm}
\norm{\nabla\widehat F_{m,U}(x)}^2
=
\bigl(h_m'(r)\beta+\eta r\bigr)^2
+a^2\bigl(\norm{g}^2-\beta^2\bigr).
\end{equation}
Applying~\eqref{eq:endpoint-H-bounds}--\eqref{eq:endpoint-B-bounds} to $F_N$ at $y$ also gives
\begin{equation}
\label{eq:endpoint-F-radial}
\ip{g}{y}
\ge
\frac12\norm{y}^2-4\sqrt N\norm{y}-\frac{13}{2}N.
\end{equation}

We now lower-bound the gradient in two cases. If $\beta\ge-\eta r/2$, then $0<h_m'(r)\le1$ implies
\[
h_m'(r)\beta+\eta r\ge\frac{\eta r}{2}.
\]
If instead $\beta<-\eta r/2$, then $\ip{g}{y}=h\beta<0$. Equation~\eqref{eq:endpoint-F-radial} therefore forces
\[
\norm{y}<(4+\sqrt{29})\sqrt N.
\]
Because $h\ge R_m=8\sqrt2\sqrt N$, we obtain
\[
\norm{z}\le s_0:=\frac{4+\sqrt{29}}{8\sqrt2}<1,
\qquad
s_0^2=\frac{45+8\sqrt{29}}{128}<\frac{7}{10}.
\]
Here $z\ne0$, since otherwise $\beta=0$. Cauchy--Schwarz then yields
\[
\norm{g}^2-\beta^2
\ge
\beta^2\left(\frac{1}{\norm{z}^2}-1\right)
\ge
\frac{3}{7}\beta^2.
\]
Moreover, $a|\beta|>\eta h/2\ge\eta R_m/2$. Thus the tangential component in
\eqref{eq:endpoint-exterior-gradient-norm} is bounded away from zero even if the radial component nearly cancels.

We finish by separating moderate and large radii. Set $K:=16$. If $R_m<r\le KR_m$, the two cases above, together
with $R_m^2=128N$, give
\begin{equation}
\label{eq:endpoint-annulus-gradient}
\norm{\nabla\widehat F_{m,U}(x)}^2
\ge
\frac{3}{7}\frac{\eta^2R_m^2}{4}
=\frac{6}{175}N.
\end{equation}
On the other hand, $\norm{y}\le h<2R_m$, $F_N(y)\le\norm{y}^2/2+3N$, and
$\widehat F_{m,U,\star}\ge0$, so
\[
\widehat F_{m,U}(x)-\widehat F_{m,U,\star}
\le
2R_m^2+3N+\frac{\eta}{2}K^2R_m^2
=\frac{5391}{5}N.
\]
Together with~\eqref{eq:endpoint-annulus-gradient}, this gives
\[
\widehat F_{m,U}(x)-\widehat F_{m,U,\star}
\le
31{,}448\norm{\nabla\widehat F_{m,U}(x)}^2.
\]

Finally, suppose $r>KR_m$ and write $s:=r/R_m\ge16$. The estimate proved in Step~2 gives
$\norm{g}\le3R_m$. The radial and tangential eigenvalues of
$J\rho_m(x)$ also give
\[
\norm{J\rho_m(x)}_{\mathrm{op}}
\le
\max\left\{e^{-(s-1)},\frac{2}{s}\right\}
\le\frac{2}{s},
\]
where $e^{s-1}\ge s$ for $s\ge1$. Since $s^2\ge16^2>240$,
\[
\norm{J\rho_m(x)^\top Ug}
\le
\frac{6R_m}{s}
\le
\frac{r}{40}
=\frac{\eta r}{2}.
\]
The quadratic regularizer therefore dominates:
\[
\norm{\nabla\widehat F_{m,U}(x)}\ge\frac{\eta r}{2}=\frac{r}{40}.
\]
Furthermore,
\[
\widehat F_{m,U}(x)-\widehat F_{m,U,\star}
\le259N+\frac{r^2}{40}
\le\frac{259}{16^2\cdot128}r^2+\frac{r^2}{40}
\le53\norm{\nabla\widehat F_{m,U}(x)}^2,
\]
where $r^2\ge16^2R_m^2=16^2\cdot128N$ was used in the second inequality.

The clipping ball, the annulus, and the large-radius region cover $\R^d$. Their bounds are all at most
$28{,}778m\norm{\nabla\widehat F_{m,U}(x)}^2$ because $m\ge m_0=5$. This proves the claimed global P\L{} bound.
\end{proof}

\end{document}
